\section{Failure window construction}

This chapter records the intended construction of a failure window
from a positive cycle.  The construction is the mathematical content
of \texttt{failureWindowExistence}; in the Lean repository that
statement remains pending, while every named ingredient below is
either compiled or compiled as a conditional interface.

\subsection{Input: a positive cycle}

Assume
\begin{equation*}
  \OrbitCycle(7).
\end{equation*}
The compiled cycle-word layer produces a starting depth $c$, a period
$p$, a state $m=\Orbit_7(c)$, and a word $w$ of length $p$ such that
\begin{itemize}[leftmargin=*]
\item $w$ is valid from $m$;
\item $\wordOrbit(w,m)=m$;
\item every entry of $w$ is the exact step weight at the
  corresponding orbit position;
\item the cycle equation
  \begin{equation*}
    2^{\wordWeight(w)}\,m=5^p\,m+\wordA(w)
  \end{equation*}
  holds.
\end{itemize}
This is the compiled input
\texttt{orbit\_\allowbreak cycle\_\allowbreak imp\_\allowbreak
cycle\_\allowbreak qb8\_\allowbreak input}.

\subsection{The QB-8 split}

The word is filtered into rise entries and C3 entries:
\begin{equation*}
  \mathrm{rise}=\{t\in w:t<3\},
  \qquad
  \mathrm{c3}=\{t\in w:t\ge3\}.
\end{equation*}
Both lists are nonempty:
\begin{itemize}[leftmargin=*]
\item a cycle cannot consist only of rise steps, because their total
  weight is at most $2P$, contradicting the cycle equation;
\item a cycle cannot consist only of C3 steps, because every C3 step
  satisfies $(5x+1)/2^t<x$ for $x\ge1$, so the states would strictly
  decrease and could not close.
\end{itemize}
The compiled lemmas
\texttt{cycleQb8Input\_\allowbreak exists\_\allowbreak rise\_\allowbreak
index} and
\texttt{cycleQb8Input\_\allowbreak exists\_\allowbreak c3\_\allowbreak
index} provide genuine positions inside the word.

\subsection{The reset start}

The construction chooses a global-minimum rise start.  The compiled
extraction lemmas
\begin{itemize}[leftmargin=*]
\item \texttt{orbit\_\allowbreak cycle\_\allowbreak imp\_\allowbreak
  min\_\allowbreak rise\_\allowbreak start};
\item \texttt{orbit\_\allowbreak reset\_\allowbreak head\_\allowbreak
  eq\_\allowbreak of\_\allowbreak rise\_\allowbreak start};
\item \texttt{orbit\_\allowbreak cycle\_\allowbreak imp\_\allowbreak
  min\_\allowbreak rise\_\allowbreak witness};
\end{itemize}
produce parameters $j,k_0,t,\delta,s$ and states $r,r_j$ such that
\begin{itemize}[leftmargin=*]
\item $t\in\{1,2\}$ and $\delta=1$ for $t=1$,
  $\delta\in\{1,3\}$ for $t=2$;
\item $s5^{k_0}=r+1$;
\item $\GeneralOrbit(r)$ and $\FullOrbit(r_j)$;
\item the reset equation
  $\mathrm{ResetHeadEq}(s,j,k_0,t,\delta,r_j)$ holds;
\item $s<5^{j-1-k_0}$, $s$ is odd, and $5\nmid s$.
\end{itemize}
This is exactly the corrected reachability predicate
$\ResetWindow(j,k_0,t,\delta,s)$.

\subsection{The block decomposition}

The block head $r_j$ is followed by a tail of length $L$ ending at the
terminal state $r_s$.  Let
\begin{equation*}
  n=s-j,\qquad \Delta=W_s-W_j,
\end{equation*}
and let $B$ be the block molecule of the tail.  The tail is a valid
word, so the exact word-orbit identity gives
\begin{equation*}
  2^{\Delta}\,r_s=5^n\,r_{j,0}+B.
\end{equation*}
The unified core studies exactly this block: the CRT candidate, the
terminal residue, and the size conditions live in the coordinates of
this tail.

\subsection{The unified-core projection}

The corrected decisive window bound is a projection of the unified
core.  In the Lean repository the interface
\begin{verbatim}
decisiveWindowValuationBoundCorrected_of_unified_core :
  unifiedCoreClosed -> decisiveWindowValuationBoundCorrected
\end{verbatim}
is compiled as a conditional interface.  Its input
\texttt{unifiedCoreClosed} is the mathematical statement whose Lean
representative is
\texttt{UnifiedCoreAudit.unified\_\allowbreak core\_\allowbreak
final\_\allowbreak no\_\allowbreak hge}; that statement is pending.

\subsection{The block-layer lower bound}

The compiled lemma \texttt{rise\_\allowbreak block\_\allowbreak
balance} gives the valuation lower bound from the reset block:
\begin{equation*}
  t=2:\qquad
  \vtwo{5^{k_0+1}s+\delta5^j}\ge2j+11,
\end{equation*}
\begin{equation*}
  t=1:\qquad
  \vtwo{5^{k_0+1}s+5^j-2}\ge2j+12.
\end{equation*}
The lower bound comes from the exact block-layer count of the rise and
C3 steps, not from a size estimate.

\subsection{The corrected upper bound}

For a reachable reset window, the corrected decisive window bound
gives
\begin{equation*}
  t=2:\qquad
  \vtwo{5^{k_0+1}s+\delta5^j}\le2j+10,
\end{equation*}
\begin{equation*}
  t=1:\qquad
  \vtwo{5^{k_0+1}s+5^j-2}\le2j+11.
\end{equation*}
The two bounds contradict the lower bounds above.  The contradiction
is packaged by the compiled lemma
\texttt{failure\_\allowbreak window\_\allowbreak contradicts\_\allowbreak
window\_\allowbreak corrected}.

\subsection{The failure window record}

The construction assembles the fields of
\begin{equation*}
  \mathrm{FailureWindow}(m,S,P,w,\mathrm{rise},\mathrm{c3},
  j,k_0,t,\delta,s):
\end{equation*}
\begin{itemize}[leftmargin=*]
\item \texttt{hj\_lt}: $j<P$;
\item \texttt{ht}: $t\in\{1,2\}$;
\item \texttt{hdelta}: the matching $\delta$;
\item \texttt{hreset}: the exact reset equation with the prefix state
  $\wordOrbit(w{\upharpoonright}j,m)$;
\item \texttt{hreach}: $\ResetWindow(j,k_0,t,\delta,s)$;
\item \texttt{hs\_odd}, \texttt{hs\_not\_five}, \texttt{hk},
  \texttt{hs\_lt}: parity, coprimality, and size;
\item \texttt{hfail\_t1}, \texttt{hfail\_t2}: the valuation lower
  bounds.
\end{itemize}
The existence of this record is
\texttt{failureWindowExistence}, pending.

\subsection{The full chain}

\begin{verbatim}
OrbitCycle 7
  -> cycleQb8Input (compiled)
  -> rise and C3 indices (compiled)
  -> reset start witness (compiled)
  -> reset equation and reachability (compiled)
  -> block tail and unified core (unifiedCoreClosed pending)
  -> decisive window bound (compiled conditional)
  -> rise_block_balance lower bound (compiled)
  -> FailureWindow fields (assembly pending)
  -> failure_window_contradicts_window_corrected (compiled)
  -> not OrbitCycle 7 (compiled conditional)
  -> IsUnboundedOrbit 7 (compiled conditional)
\end{verbatim}

\subsection{What is and is not claimed}

Every step before the failure-window assembly is either compiled or a
direct projection of the unified core.  The paper does not claim that
\texttt{failureWindowExistence} is compiled: it is the pending input
of the cycle bridge.  The construction also does not claim that the
unified core is closed; that closure is exactly the pending statement
\texttt{unified\_\allowbreak core\_\allowbreak final\_\allowbreak
no\_\allowbreak hge}.

\subsection{Status ledger}

\begin{center}
\small
\begin{tabular}{@{}p{0.60\textwidth}p{0.34\textwidth}@{}}
\toprule
Statement & Status \\
\midrule
cycle-word extraction and QB-8 split & compiled \\
rise-start reset witness & compiled \\
block-layer lower bound & compiled \\
failure-window contradiction lemma & compiled \\
corrected bound as a projection & compiled (conditional) \\
\texttt{failureWindowExistence} & pending \\
\texttt{unified\_\allowbreak core\_\allowbreak final\_\allowbreak
no\_\allowbreak hge} & pending \\
\bottomrule
\end{tabular}
\end{center}
